\documentclass[a4paper,11pt]{ltjsarticle}
\usepackage{amsmath,amssymb}
\usepackage{enumitem}
\usepackage{geometry}
\geometry{margin=25mm}

\title{党勢の趨勢を指数重み付き局所2次回帰で推定する案}
\author{}
\date{2026年10月}

\begin{document}
\maketitle

本稿では、党勢指数$C_p(T)$の前月差に代わる「直近の趨勢」の指標として、指数重み付きの局所2次回帰による傾きの推定を提案する。まず現行の前月差の問題を確認する(\S\ref{sec:problem})。次に記号を定義し、回帰モデルと推定量を示す(\S\ref{sec:notation}--\ref{sec:estimator})。続いて現行指数との関係(\S\ref{sec:relation})、誤差の評価と最小二乗法の前提(\S\ref{sec:inference})、傾きに寄与する選挙と寄与しない選挙(\S\ref{sec:contribution})を述べ、最後に検証で決める事項をまとめる(\S\ref{sec:validation})。

\section{現行の前月差の問題}
\label{sec:problem}

現行の前月差は、その月の党勢の変化だけでなく、12か月前の選挙の影響も同じ大きさで含んでしまう。$C_p(T)$は直近12か月の窓に投票日があった選挙を一括して集計した値である。したがって、窓を1か月進めたときの差は、新しく窓に入った選挙の効果と、窓から外れた選挙の効果の和になる。第1回(2026年9月)では、前月差の大半が後者、すなわち2025年9月の知事選挙が窓から外れた効果だった。直近の趨勢を見たいという目的からすると、1年前の選挙が抜けたことで値が動くのは望ましくない。

窓をやめて各月の選挙だけで値を作り、その系列を差分する方法も考えられるが、これは別の問題を生む。月によって含まれる階層が大きく異なるからである。例えば2026年2月はほぼ衆院選だけで値が決まり、多くの月は市区町村議会と市区町村長だけで決まる。政党の占有率の水準は階層ごとに大きく違うので、この系列は党勢と無関係に、階層の構成が変わるたびに跳ぶ。以下の方法は、この2つの問題を同時に避けることを狙う。

\section{記号}
\label{sec:notation}

本節では、回帰モデルを書くのに必要な記号を定義する。選挙イベント$e$、その代表人口$P_e$、議席数$n_e$、政党$p$の占有率$\phi_p(e)$、階層$\ell\in L$は設計文書(\texttt{design\_document.tex}、\S2)と同じ意味で使う。ただし設計文書の$T$は「期間」を表すので、本稿では評価する時点を別の記号$\tau$で表す。新しく導入する記号は次のとおりである。

\begin{itemize}
  \item $\tau$: 評価時点(各月の月末)。
  \item $t_e$: 選挙イベント$e$の投票日。
  \item $h$: 記憶の長さ(月)。時間方向の重みがどれだけ速く減衰するかを決める。
  \item $s_e$: 評価時点から見た$e$の相対時刻で、$s_e=(t_e-\tau)/h$と定める。$t_e\le\tau$の選挙では$s_e\le 0$であり、$s_e=0$が評価時点に当たる。
  \item $K(s)$: 時間方向の重み。指数減衰$K(s)=e^{s}$を使う。$h=12$か月のとき、重みが半分になるのは約8か月前の選挙である。
  \item $E(\tau)$: $t_e\le\tau$を満たす全ての選挙イベント$e$の集合。期間で区切らず、評価時点以前の全データを使う。
  \item $E_p(\tau)$: $E(\tau)$のうち、政党$p$の結党日以降に投票日がある選挙イベントの集合。結党前の選挙は、その党の回帰には使わない。
  \item $\tilde\phi_p(e)$: 無所属を除いて正規化した占有率。$\phi_{\text{無所属}}(e)<1$の選挙について$\tilde\phi_p(e)=\phi_p(e)/(1-\phi_{\text{無所属}}(e))$と定める。当選者が全員無所属の選挙は、この値が定義できないので回帰から除く。
  \item $\ell(e)$: 選挙イベント$e$が属する階層。ただし市区町村長は、推薦・支持の名簿で仕分け済みの選挙と未仕分けの選挙を別の階層として扱う(理由は\S\ref{sec:contribution})。この拡張した階層の集合も$L$と書く。
\end{itemize}

\section{回帰モデル}
\label{sec:model}

提案するモデルは、各階層に固有の水準を許しつつ、時間の変化は全階層で共通の2次式で表す、というものである。回帰の1点は月ごとの集計値ではなく選挙イベント1件である。政党$p$について、$E_p(\tau)$に属する各選挙イベント$e$の占有率を次の式で表す。
\begin{equation}
  \tilde\phi_p(e) = \alpha_{p,\ell(e)} + \beta_p\, s_e + \gamma_p\, s_e^2 + \varepsilon_e。
  \label{eq:model}
\end{equation}
各パラメータの意味は次のとおりである。

\begin{itemize}
  \item $\alpha_{p,\ell}$: 階層$\ell$における政党$p$の占有率の水準(評価時点$s=0$での値)。階層ごとに別の値を持つ。
  \item $\beta_p$: 評価時点での傾き。全階層で共通とする。
  \item $\gamma_p$: 曲率。全階層で共通とする。
  \item $\varepsilon_e$: 個々の選挙のばらつき(誤差項)。
\end{itemize}

階層ごとに切片$\alpha_{p,\ell}$を分けるのが、このモデルの要点である。例えば、ある政党の占有率が市区町村長では0.5、市区町村議会では0.3で、どちらも時間的に一定だとする。前半に市区町村長選挙が多く、後半に市区町村議会選挙が多いと、$\tilde\phi_p(e)$を時刻に対して素朴に並べた点列は右下がりに見える。しかし式\eqref{eq:model}では0.5と0.3の差を$\alpha_{p,\ell}$が吸収するので、$\beta_p=\gamma_p=0$と推定される。つまり、階層の構成が時期によって偏っていても、それが傾きに混入しない。

1次式ではなく2次式にするのは、評価時点での傾きを偏りなく推定するためである。1次式で当てはめると、$\beta_p$は記憶の長さ全体を均した平均的な傾きになり、趨勢が途中で反転した場合には評価時点の傾きとずれる。2次の項$\gamma_p s_e^2$を加えれば、曲がった趨勢にも追従できる。$s_e^2$の微分は$s=0$で0になるので、評価時点の傾きは$\gamma_p$によらず$\beta_p$だけで決まる。

3次以上にしないのは、端点での傾きの分散が次数とともに急に増えるためである。選挙が時間的に一様に分布し、切片が1つだけの単純な場合に、$K(s)=e^s$のもとで$\hat\beta$の分散を計算すると、1次式を1として2次式は5、3次式は14になる。後述の$D_p$の分散は$h^{-3}$に比例するので、3次式で2次式と同じ分散を得るには$h$を約1.4倍に延ばす必要がある。党勢の変化は新党の急伸や合流のような段差に近い形をとることが多く、高い次数で拾える滑らかな構造は乏しいと考えられるので、2次式で十分とした。

\section{推定量}
\label{sec:estimator}

パラメータは重み付き最小二乗法(WLS)で推定する。各選挙イベントの重みは、時間方向の重み$K(s_e)$と、現行指数での階層内・階層間の重みの積とする。現行指数の$N_\ell(T)$・$\Pi_\ell(T)$は窓の中の単純な合計だったが、本案では$K$で重み付けした合計に置き換える。
\begin{gather}
  N_\ell(\tau) = \sum_{e\in E(\tau),\,\ell(e)=\ell} K(s_e)\,n_eP_e,\qquad
  W_\ell(\tau) = \sqrt{N_\ell(\tau)},\qquad
  \Pi_\ell(\tau) = \sum_{e\in E(\tau),\,\ell(e)=\ell} K(s_e)\,P_e,\\
  \omega_e = K(s_e)\cdot
  \frac{W_{\ell(e)}(\tau)}{\sum_{\ell'\in L} W_{\ell'}(\tau)}\cdot
  \frac{P_e}{\Pi_{\ell(e)}(\tau)}。
  \label{eq:weight}
\end{gather}
推定量は次の最小化問題の解として定める。
\begin{equation}
  (\hat\alpha_p, \hat\beta_p, \hat\gamma_p)
  = \operatorname*{arg\,min}_{\alpha,\beta,\gamma}
  \sum_{e\in E_p(\tau)} \omega_e
  \bigl(\tilde\phi_p(e) - \alpha_{\ell(e)} - \beta s_e - \gamma s_e^2\bigr)^2。
  \label{eq:wls}
\end{equation}
ここで$\hat\alpha_p=(\hat\alpha_{p,\ell})_{\ell\in L}$である。$s_e$は$h$で割って無次元化しているので、$\hat\beta_p$は「$h$か月あたりの占有率の変化」を表す。記事で報告する月あたりの趨勢は
\begin{equation}
  D_p(\tau) = \frac{\hat\beta_p}{h}
\end{equation}
とする。これは、当てはめた曲線$\hat\alpha_{p,\ell}+\hat\beta_p s+\hat\gamma_p s^2$を実時間$t$で微分し、評価時点$t=\tau$で評価した値に等しい。

推定量は$\tilde\phi_p(e)$について線形であり、次の2つの性質を持つ。

\begin{enumerate}
  \item 全政党の傾きの合計は0になる。説明変数と重みは政党によらないので、$E_p(\tau)$が全政党で同じなら、$\sum_p\tilde\phi_p(e)=1$から$\sum_p D_p(\tau)=0$が従う。ある政党の趨勢が上向きなら、他のどこかが同じだけ下向きになる。結党前の選挙を除く政党がある期間は、この等式は近似的にしか成り立たない。
  \item 傾きを選挙ごとの寄与に分解できる。$\hat\beta_p=\sum_{e}c_e\,\tilde\phi_p(e)$と書ける。係数$c_e$は説明変数と重みだけで決まる。したがって$D_p(\tau)$のうち、どの選挙がどれだけ寄与したかを正確に示せる。現行の記事の「窓に入った効果・外れた効果」の分解は、この寄与の分解で置き換える。
\end{enumerate}

\section{現行指数との関係}
\label{sec:relation}

提案する推定量は、現行の$C_p(T)$を特別な場合として含む。$K$を直近12か月で1、それより前で0とする階段関数に置き換え、$\beta=\gamma=0$に固定し、正規化前の$\phi_p(e)$を使って式\eqref{eq:wls}を解く。このとき各階層の切片は階層内の$P_e$加重平均$\hat\alpha_{p,\ell}=\phi_p^\ell(T)$になる。これを$W_\ell$で合成した値$\sum_\ell W_\ell\hat\alpha_{p,\ell}/\sum_\ell W_\ell$は、現行の$C_p(T)$そのものである。つまり、現行指数は「12か月の中で党勢は一定」と仮定した当てはめであり、本案はその仮定を外して時間変化を許したものと位置づけられる。

指数減衰の重みを使うことで、窓から外れる選挙による段差はなくなる。現行の$C_p(T)$では、12か月前の選挙が直近の選挙と同じ重みを持ったまま、ある月に突然抜ける。本案では古い選挙の重みが連続的に減っていくので、ある月に突然データが抜けることがない。$C_p(T)$は水準の参考として残し、記事の主役を$D_p(\tau)$に移す。

\section{誤差の評価と最小二乗法の前提}
\label{sec:inference}

本案で最小二乗法の前提として確かめる必要があるのは、平均の構造が正しいかどうかと、誤差が大きな塊ごとに独立かどうかの2点である。古典的な最小二乗法の前提には、誤差の分散が一定であること、誤差どうしが無相関であること、誤差が正規分布に従うことも含まれる。しかし分散一定と無相関は、推定量が最も効率的であるための条件と、古典的な標準誤差の公式が正しいための条件にすぎない。本案では標準誤差を投票日ごとのクラスタに頑健な形で計算するので、この2つは要らない。正規性も、クラスタが十分多ければ中心極限定理で不要になる。

残る2つの前提は、具体的には次のように確認する。

\begin{enumerate}
  \item 平均の構造。$\hat\beta_p$が評価時点の傾きを偏りなく推定するには、式\eqref{eq:model}の誤差$\varepsilon_e$の平均が$s_e$によらず0である必要がある。2次式の当否は\S\ref{sec:validation}の予測誤差で確かめる。より深刻なのは、どの自治体で選挙があったかが時期によって偏る場合である。例えば統一地方選挙で改選される自治体と、それ以外の時期に改選される自治体とで、政党の地盤が系統的に違えば、その差が傾きに混入する。統一地方選挙の選挙とそれ以外の選挙とで残差の平均を比べて確認する。
  \item クラスタ間の独立性。同じ投票日の選挙は全国的な情勢を共有するので、誤差は投票日ごとに相関しうる。クラスタに頑健な標準誤差はこの相関を許すが、異なる投票日の間の独立性は仮定する。隣り合う投票日の残差に相関があれば、クラスタを月単位に広げる。また指数重みのもとでは実効的なクラスタ数が少なくなるので、少数クラスタの補正が要るかどうかも確かめる。
\end{enumerate}

誤差の分布の裾が重く、分散が発散することはない。$\tilde\phi_p(e)$は$[0,1]$に収まるので、誤差$\varepsilon_e$も有界であり、全てのモーメントが有限だからである。裾の重さが問題になるのは、誤差ではなく重みの側である。重みは$P_e$に比例するが、市区町村の人口は上位でべき分布に近い。手元の人口データ(1709自治体)で上位$k$自治体からHill推定量を計算すると、裾の指数は$k=50$で1.7、$k=400$で1.1であり、いずれも2を下回る。指数が2未満のべき分布は分散が発散する。

この裾の重さは、推定の実効的な標本の大きさを大きく減らす。重みを固定すれば誤差の分散は有限なので、推定量の分散は$\sigma^2/n_{\text{eff}}$の形で書ける。ここで$\sigma^2$は誤差の分散、$n_{\text{eff}}=(\sum_e\omega_e)^2/\sum_e\omega_e^2$は実効標本数である。全自治体を人口で重み付けすると$n_{\text{eff}}\approx 218$になる。1年分の改選に相当する430自治体を無作為に選ぶと$n_{\text{eff}}$の中央値は61で、最大の自治体1つが重みの7\%(95パーセンタイルで12\%)を占める。さらに指数重みをかけると$n_{\text{eff}}$はもっと小さくなる。また裾の指数を$a$($1<a<2$)とすると、$n_{\text{eff}}$は選挙数$n$に対しておおむね$n^{2-2/a}$でしか増えない。記憶の長さ$h$を延ばしても、精度は選挙数ほどには改善しない。

そこで、次の3つの対策をとる。

\begin{enumerate}
  \item $D_p(\tau)$と一緒に$n_{\text{eff}}$と、重みが最大の選挙の割合を記録する。
  \item 標準誤差の正規近似は、少数の大都市が重みを占めると不正確になる。そのため、投票日単位のワイルド・クラスタ・ブートストラップによる信頼区間と比べて確認する。
  \item 各選挙の寄与$c_e\,\tilde\phi_p(e)$を見て、横浜市長選挙や東京都知事選挙のような1件の選挙が$D_p$を決めていないかを確認する。
\end{enumerate}

階層内の重みを$P_e$から$\sqrt{P_e}$のように圧縮すれば裾は軽くなる。しかしこれは階層内の集計方法という指数の定義そのものの変更であり、本案の範囲外とする。

重みの意味についても補足する。$\omega_e$は、各選挙が人口をどれだけ代表するかという指数の定義から決めたもので、誤差の分散から決めたものではない。誤差の相関と分散を推定して重みに反映する一般化最小二乗法(GLS、ここでは投票日のランダム効果を持つ混合モデル)を使えば、推定は効率的になる。ただし傾きが自治体ごとに異なる場合、重みが変わると「どの平均の傾きを推定しているか」も変わる。そのため本推定は人口の重みによるWLSとし、GLSは感度分析として並べる。

\section{傾きに寄与する選挙と寄与しない選挙}
\label{sec:contribution}

専用の切片を持つ階層の選挙は、その階層の中に投票日の異なる選挙が複数ある場合に限って傾きの推定に寄与する。極端な場合として、ある階層に選挙が1件しかないとする。この1点には専用の切片$\alpha_{p,\ell}$があり、$\beta_p,\gamma_p$をどう選んでも$\alpha_{p,\ell}$を調整すれば残差を0にできる。したがってこの点を除いて推定しても$\hat\beta_p,\hat\gamma_p$は変わらない。$\beta_p$が全階層共通であることは、この結論を変えない。階層の水準を自由にした以上、ある階層の選挙を別の階層の水準と比べる手段がないからである。

この性質から、国政選挙の扱いについて次のことが言える。衆院選・参院選はそれぞれ1回の選挙が1つのイベントであり、別々の階層として扱う。全期間を使うので、各院の階層には過去の選挙も含まれ、それらとの比較を通じて国政選挙も傾きに寄与する。その寄与の大きさは、2回目以降の選挙の重み$K(s_e)$で決まる。例えば$h=12$か月の場合、2026年10月時点で2024年10月の衆院選の重みは$e^{-2}\approx0.14$である。

名義の変更も傾きに現れる。立憲民主党は、2024年の衆院選では議席を得ているが、2026年の衆院選では議席が中道改革連合の名義になっており0である。この差は立憲民主党の下降の傾きとして$\hat\beta_p$に現れる。本プロジェクトでは窓の段差を党勢の低下と読む方針をとっており、名義の変更による傾きも同様に、その名義での党勢の低下として扱う。記事ではこの点を説明する。

同じ性質を、市区町村長の推薦・支持データの遅れへの対策に使う。名簿の年版は収録期間の終わりから9〜10か月後に公開されるので、直近およそ1年半の市区町村長選挙は、当選者が無所属のまま推薦・支持政党に仕分けられない。この期間の選挙は無所属の割合が高く、評価時点に近いほど政党の値が下がって見える。これをそのまま回帰すると、傾きが系統的に下向きにずれる。そこで未仕分けの選挙を別の階層として専用の切片を与える(\S\ref{sec:notation})。こうすると、仕分け済みの選挙との水準差は切片が吸収し、傾きには未仕分けの選挙どうしの時間変化だけが効く。新しい年版が公開されたら、仕分けられた選挙を仕分け済みの階層に移して過去の月も再計算する。

\section{検証で決める事項}
\label{sec:validation}

本案で未定のまま残るのは、記憶の長さ$h$だけである。$h$は、評価時点の当てはめ曲線を少し先まで延ばしたときの予測の当たり具合で選ぶ。具体的には、2020年以降の各月末$\tau$について、時点$\tau$での推定値$\hat\alpha_{p,\ell}+\hat\beta_p s+\hat\gamma_p s^2$で、その後3か月の選挙の$\tilde\phi_p(e)$を予測し、重み付きの二乗誤差を集計する。候補は$h=6,9,12,18$か月とする。真の党勢の趨勢は観測できないが、この基準なら一方の側のデータだけを使う推定量の性能を直接比べられる。

データは2018年から存在するので、最初の$2h$か月程度は記憶が途中で切れている。この立ち上がりの期間の値は公表しない。

次の2点は本案の範囲外とし、将来の課題とする。

\begin{itemize}
  \item 自治体の固定効果。全期間を使えば同じ自治体の選挙が4年ごとに複数回含まれるので、自治体ごとの地盤の違いを切片で吸収できる。これは\S\ref{sec:inference}で述べた「どの自治体で選挙があったか」による偏りへの直接の対策になる。ただし比較が4年おきの組になるので、$h$が24か月程度以上でないと効かない。
  \item 衆院選と参院選を1つの階層にまとめる案。まとめれば国政選挙が傾きにより多く寄与するが、衆参の水準差が一定だという仮定が入る。
\end{itemize}

\end{document}
